\documentclass[11pt]{article}

\usepackage[left=1.25in,right=1.25in,top=1in,bottom=1in]{geometry}
\usepackage{setspace}
\usepackage{microtype}
\usepackage{amsmath,amssymb,amsfonts}
\usepackage{mathtools}
\usepackage{graphicx}
\usepackage{hyperref}

\setstretch{1.15}
\setlength{\parindent}{0pt}
\setlength{\parskip}{0.75em}

\title{The Politics of Oracle-Free Evaluation:\\
Constraint, Generativity, and the Problem of Validation}

\author{Flyxion}
\date{\today}

\begin{document}

\maketitle

\begin{abstract}

Recent proposals for ``oracle-free'' evaluation of mathematical reasoning, most prominently Consequence-Based Utility (CBU), seek to validate candidate solutions without access to ground-truth answers by measuring their ability to generate correct reasoning across curated neighborhoods of related problems (Son et al., 2026). These approaches replace verification by correspondence with verification by generativity. While technically innovative, this substitution encodes substantive epistemological commitments. It presumes that structural fertility within a formal manifold can function as a proxy for correctness, and that validation may be relocated from shared disciplinary practices into computational infrastructures.

This essay argues that consequence-based validation is reliable only within closed formal domains and becomes systematically biased when extended to open, physically instantiated, or temporally evolving systems. I introduce the notion of \emph{constraint-indexed correctness}, according to which the validity of an explanation is determined not merely by inferential propagation but by its stability under transformations imposed by physical law, scale transition, measurement, and historical evolution. The analysis situates CBU within a longer history of shifts from correspondence to operational adequacy (Kuhn, 1962; Lakatos, 1976; Hacking, 1983) and examines how contemporary platform-centered infrastructures reshape what counts as knowledge (Lessig, 1999; Marres, 2017).

Rather than rejecting consequence-based reasoning outright, the paper proposes a hybrid research program in which generative validation operates as an inner-loop heuristic embedded within outer-loop constraint tests. Such an approach restores what earlier philosophies of science treated as indispensable: friction with reality as the ultimate arbiter of explanation.

\end{abstract}

\newpage
\section{Introduction}

The aspiration to evaluate reasoning without recourse to authoritative answers is not new. Mathematics has long advanced through conjectural exploration in regions where verification is costly or delayed. What is new is the attempt to mechanize this process at scale. Consequence-Based Utility proposes that the correctness of a candidate solution can be inferred from its downstream usefulness: a solution is judged by how effectively it enables a model to solve related problems (Son et al., 2026).

This proposal appears at first glance to mirror ordinary mathematical practice. Mathematicians often gain confidence in an argument because it organizes a field, generates corollaries, or simplifies adjacent results. Yet historical studies of mathematics emphasize that such confidence emerges through prolonged communal scrutiny, reconstruction, and contestation (Lakatos, 1976). Proof is not merely written but socially stabilized. Scientific knowledge more broadly relies upon norms of reproducibility, organized skepticism, and shared validation (Merton, 1973; Shapin, 1994).

CBU replaces this distributed process with a computational analogue. Instead of adversarial uptake over time, it evaluates candidate reasoning through short cycles of neighborhood perturbation and automated scoring. Validation becomes a question of statistical transfer rather than reconstructive understanding.

This shift reflects a broader transition in twentieth- and twenty-first-century science from correspondence theories of truth toward operational adequacy. Kuhn's analysis of paradigm-dependent evaluation (Kuhn, 1962) and Hacking's emphasis on intervention rather than representation (Hacking, 1983) both describe forms of validation grounded in practice. However, those accounts retained an essential feature absent in oracle-free evaluation: sustained interaction with a resistant world.

The central question of this essay is therefore not whether consequence-based evaluation functions technically, but what kind of correctness it is capable of detecting.

\section{From Correspondence to Consequence}

Classical accounts of knowledge distinguished between two intertwined criteria. A claim was considered valid if it corresponded to invariant structure and if it integrated coherently into a broader inferential system. In mathematics these criteria often coincide because the domain itself is formal. In empirical sciences they diverge.

Operational approaches risk collapsing the distinction. When success is measured by predictive or inferential performance alone, correctness becomes indistinguishable from calibration. Goodhart's observation that optimized proxies diverge from intended targets (Goodhart, 1984) applies as readily to epistemology as to economics.

CBU embodies this tension. By evaluating candidate reasoning through induced performance on related tasks, it privileges transportability over constraint. The resulting signal measures what might be called \emph{intra-model coherence}. It does not test whether the reasoning remains stable under transformations imposed by measurement, scale, or embodiment.

Philosophers of science have long warned that explanatory adequacy cannot be reduced to predictive success alone. Cartwright argued that laws function only within constrained regimes rather than universally (Cartwright, 1983), while Anderson showed that behavior changes qualitatively across scales (Anderson, 1972). These insights imply that validation must attend to the conditions under which reasoning remains invariant, not merely to the consequences it produces.

\section{The Operational Structure of Consequence-Based Utility}

CBU replaces an unavailable correctness oracle with a performance functional. Let \(Q\) denote a target problem, \(C\) a candidate solution, and \(\mathcal{N}(Q)\) a curated set of related questions. A solver model conditioned on \(C\) is evaluated on \(\mathcal{N}(Q)\), and the resulting accuracy score is treated as evidence for or against the validity of \(C\) (Son et al., 2026).

The method is ingenious in that it transforms evaluation into a measurable quantity even in the absence of known answers. Yet the procedure presupposes that correctness is equivalent to transferable method-level information. This assumption holds naturally in closed symbolic systems. It is far less secure in domains where explanations must survive scale transitions, causal interventions, or measurement constraints.

The remainder of this paper develops a framework for understanding these limitations and for extending evaluation beyond consequence alone.

\section{Constraint-Indexed Correctness}

To move beyond purely consequence-based validation, we require a notion of correctness that is sensitive to the conditions under which an explanation remains valid. I introduce here the concept of \emph{constraint-indexed correctness}, according to which the validity of a candidate explanation is defined relative to a structured set of admissible transformations rather than to a neighborhood of formally similar problems.

Let
\[
\mathcal{K} = (\Phi, \Sigma, \Tau, \Mu)
\]
denote a constraint tuple consisting of four components. The first, \(\Phi\), represents admissible physical or logical invariants, such as conservation relations or definitional identities. The second, \(\Sigma\), denotes scale transformations mapping descriptions across levels of resolution. The third, \(\Tau\), encodes temporal evolution operators describing how systems change under dynamical progression. The fourth, \(\Mu\), represents measurement interfaces that specify how theoretical quantities are operationally accessed.

Given a candidate explanation \(C\) addressing a question \(Q\), we say that \(C\) is \emph{correct relative to} \(\mathcal{K}\) if its predictive and inferential structure remains stable under all admissible transformations generated by \(\mathcal{K}\). Formally,
\[
C \models_{\mathcal{K}} Q
\]
when, for every transformation \(T\) induced by \(\Phi\), \(\Sigma\), \(\Tau\), or \(\Mu\), the transformed explanation \(T(C)\) continues to entail the appropriately transformed claim \(T(Q)\).

This definition shifts validation from propagation across syntactic neighbors to invariance under constraint. Where CBU asks whether a solution generalizes across related prompts, constraint-indexed validation asks whether the explanation survives contact with the structure of the world.

The distinction echoes long-standing insights in the philosophy of science. Explanations gain credibility not merely by predicting new cases but by remaining stable under interventions and regime changes (Hacking, 1983; Pearl, 2009).

\section{A Toy Separation Result}

It is possible to construct situations in which consequence-based evaluation cannot distinguish between candidates that constraint-indexed validation separates.

Consider two candidate derivations \(C_1\) and \(C_2\) that yield identical algebraic predictions across a curated neighborhood of problems. Suppose \(C_1\) respects a conserved quantity encoded in \(\Phi\), while \(C_2\) introduces a hidden violation that cancels locally but accumulates under scaling transformations \(\Sigma\).

Within a finite neighborhood, both candidates generate equivalent downstream success. CBU therefore assigns comparable utility scores. However, under repeated application of a scale operator \(S \in \Sigma\), the predictions of \(C_2\) diverge systematically, while \(C_1\) remains invariant.

Thus there exist pairs \((C_1, C_2)\) such that
\[
U_{\mathrm{CBU}}(C_1) \approx U_{\mathrm{CBU}}(C_2)
\quad \text{but} \quad
C_1 \models_{\mathcal{K}} Q \;\; \text{and} \;\; C_2 \not\models_{\mathcal{K}} Q.
\]

This separation illustrates that consequence-based validation measures local transferability, whereas constraint-indexed validation detects structural compatibility across regimes.

Analogous phenomena appear throughout science. Renormalization in physics reveals theories that agree at one scale yet diverge at another (Anderson, 1972). Biological scaling laws invalidate models that extrapolate linearly across organism size. Economic equilibria fail when institutional constraints shift (Scott, 1998). These are not edge cases but characteristic features of open systems.

\section{Temporal Functional Evaluation}

CBU treats validation as effectively static. A candidate solution is scored through a bounded sampling procedure. Yet many explanations must remain reliable across unfolding time, not merely across adjacent formal variants.

Let observational updates be indexed by \(t\). Define a temporal validation functional
\[
U_t(C) = \mathbb{E}\bigl[\mathcal{P}(C \mid \text{constraints observed up to } t)\bigr].
\]
Correctness now depends on the trajectory of constraint satisfaction rather than a single evaluation window.

Static consequence-based evaluation appears as a limiting case when temporal evolution is assumed independent of explanatory structure. Such independence rarely holds in real systems, where feedback, adaptation, and path dependence dominate (Tetlock, 2005).

The shift from static to temporal validation reintroduces what Kuhn described as the historical dimension of scientific change, in which theories persist only by surviving successive confrontations with new data (Kuhn, 1962).

\section{Embodiment and Instantiation Cost}

Formal structures are cheap. Instantiated explanations are not. Any evaluation regime that ignores embodiment risks privileging models that are internally coherent yet physically unrealizable.

Let \(\mathcal{I}\) denote an instantiation operator mapping abstract descriptions into realizable processes constrained by energy, material, and measurement limits. A candidate explanation is viable only if \(\mathcal{I}(C)\) exists and preserves predictive structure under those constraints.

This emphasis aligns with Polanyi's account of tacit knowledge and materially grounded practice (Polanyi, 1966), as well as with studies of craft and skilled engagement in technological work (Sennett, 2008; Crawford, 2009). Knowledge is not merely encoded; it is enacted within constrained environments.

CBU intentionally abstracts away from instantiation cost in order to scale. Constraint-indexed validation instead treats embodiment as a primary filter distinguishing meaningful structure from formal excess.

\section{Information, Compression, and External Reference}

The contrast between consequence and constraint can also be expressed in information-theoretic terms. CBU evaluates how well a candidate compresses reasoning within a constructed neighborhood. Constraint-indexed validation asks whether that compression preserves mutual information with an external system.

In Shannon's framework, efficient encoding does not guarantee semantic fidelity (Cover and Thomas, 2006). A model may achieve excellent compression while discarding structure necessary for intervention or extrapolation. Constraint-indexed validation therefore prioritizes preservation of externally anchored information over internal efficiency.

\section{Infrastructure and the Institutionalization of Evaluation}

Evaluation methods are not epistemically neutral; they inherit the affordances of the infrastructures that support them. Digital platforms shape participation, access, and authority through architectural design (Lessig, 1999). Studies of digital knowledge production emphasize how platforms reorganize epistemic practices around metrics, visibility, and control (Marres, 2017).

The migration of validation into model-centric ecosystems therefore alters not only how solutions are scored but who can verify those scores. Earlier scientific regimes relied upon durable archival systems and independent reproducibility (Merton, 1973). Platform-mediated validation risks replacing those norms with access-dependent evaluation pipelines.

This is not an argument against new infrastructures, but a reminder that epistemology is inseparable from its material supports.

\section{Toward Hybrid Validation Architectures}

The critique advanced here does not entail abandoning consequence-based evaluation. Transferability across related problems remains a valuable heuristic. The issue is sufficiency.

A more robust architecture would treat consequence-based scoring as an inner-loop filter identifying candidates worthy of further scrutiny. Outer loops would then test invariance under constraint transformations, temporal updates, and instantiation requirements.

Such a hierarchy resembles layered verification systems in other domains, where fast approximate checks precede slower but more reliable validation. It restores the balance between scalability and resistance that has historically characterized successful scientific methodologies.

\section{Geometry, Constraint, and Learning Dynamics}

These considerations suggest a reinterpretation of learning itself. Standard optimization treats parameter space as an undifferentiated Euclidean domain. A constraint-aware perspective instead models learning as motion restricted to an admissible manifold.

Information geometry provides a mathematical language for such constrained motion, where updates are projected onto structures preserving relevant invariants (Amari, 2016). Learning becomes alignment with structure rather than unrestricted exploration.

From this standpoint, elaborate evaluation schemes compensate for an overly permissive search space. If admissible trajectories were constrained from the outset, many formally generative but meaningless candidates would never arise.

The contrast mirrors earlier tensions between exploratory and constraint-driven scientific approaches. One expands possibility and filters afterward. The other encodes structure into the dynamics of inquiry itself.

\section{Adversarial Considerations and the Limits of the Critique}

Any critique of consequence-based validation must address its strongest defense. Within strictly formal domains, one may argue that consequence is indistinguishable from truth. If mathematical objects are defined entirely by inferential roles, then the capacity of a candidate solution to organize correct reasoning across related cases appears to capture precisely what correctness means.

There is historical precedent for this view. Structuralist interpretations of mathematics treat objects as positions in relational systems rather than as independently existing entities. From that standpoint, validation through coherent extension is not a proxy but an instantiation of mathematical truth itself.

The present argument does not deny this possibility. Instead, it asserts that the equivalence between consequence and correctness holds only under closure conditions that rarely extend beyond pure formalism. Once explanations interact with measurement, scale transitions, or embodied implementation, inferential fertility ceases to guarantee adequacy.

A second objection holds that scientific practice itself often validates theories through coherence rather than direct correspondence. Theories are rarely rejected immediately upon anomaly; they are adjusted, extended, and reinterpreted (Lakatos, 1976). Yet such flexibility operates within a landscape constrained by empirical resistance. Adjustments must eventually reconcile with observation or fail to stabilize historically.

A third objection notes that constraint-indexed validation also relies on human judgment in specifying admissible constraints. This is correct. Constraints are not discovered from nowhere. The difference lies in how judgment enters the system. In consequence-based evaluation, human expertise constructs neighborhoods whose epistemic role remains implicit. In constraint-indexed validation, the role of those assumptions is made explicit and becomes itself subject to scrutiny and revision.

The disagreement, therefore, is not between objective automation and subjective intervention, but between implicit and explicit embedding of assumptions.

\section{Computational Cost and the Scalability Objection}

A natural concern is that constraint-indexed validation may be computationally more demanding than oracle-free evaluation. Consequence-based methods scale efficiently because they reduce validation to repeated inference within a fixed model. Constraint-aware approaches require additional structure, simulation, or measurement.

This asymmetry is real and historically familiar. Empirical science has always advanced more slowly than formal speculation because contact with reality imposes cost. Instrumentation, replication, and cross-scale analysis demand resources that purely symbolic reasoning does not.

The relevant question is not whether constraint-indexed validation is cheaper, but whether its additional expense corresponds to epistemic value. If correctness in open systems requires confrontation with constraints, then any method that eliminates that confrontation risks trading reliability for throughput.

Rather than replacing consequence-based evaluation, constraint-aware approaches should be understood as selectively applied where domain structure demands them.

\section{Constraint Discovery and the Bootstrap Problem}

A further challenge concerns the origin of the constraint tuple \(\mathcal{K}\). How can validation depend on constraints that are themselves imperfectly known?

Scientific practice resolves this bootstrap problem iteratively. Constraints are neither fixed nor arbitrary; they are discovered, refined, and occasionally replaced through cycles of prediction and failure. Measurement reveals regularities, which become provisional invariants subject to further testing.

Constraint-indexed validation therefore operates not with immutable conditions but with evolving approximations. The tuple \(\mathcal{K}\) is historically extended rather than presupposed. This view aligns with accounts of scientific development emphasizing feedback between theory and instrumentation (Hacking, 1983).

\section{Worked Illustration: Local Transfer Versus Scale Stability}

To illustrate the distinction concretely, consider an explanation that models a system through linear approximation. Within a curated neighborhood of small perturbations, the approximation performs well and supports accurate reasoning. Consequence-based evaluation would assign it high utility.

However, when subjected to scale transformation, nonlinear terms become dominant, and predictions diverge. Constraint-indexed validation identifies this failure because invariance under \(\Sigma\) is violated.

Such examples appear across disciplines. Linear elasticity fails at large strain, mean-field approximations collapse near phase transitions, and rational-agent models break under institutional change. In each case, local generativity coexists with global inadequacy.

\section{Institutional Implications of Validation Regimes}

Evaluation frameworks shape not only epistemology but institutional practice. Systems that privilege rapid, automated scoring encourage research cultures oriented toward benchmark optimization. Systems that emphasize constraint confrontation reward slower, experimentally grounded inquiry.

Historical analyses show that scientific authority has often followed the tools of validation. The emergence of laboratory instrumentation reorganized epistemic credibility around reproducible experiment. Digital infrastructures risk reorganizing credibility around platform-mediated metrics unless counterbalanced by practices that preserve independent verification (Merton, 1973; Lessig, 1999).

Recognizing these dynamics is not an indictment of new methods but a call to examine how evaluative mechanisms redistribute epistemic authority.

\section{A Research Program for Constraint-Indexed Validation}

The framework proposed here suggests several avenues for future work. One direction is the construction of evaluation tasks explicitly designed to test invariance under scale or causal intervention rather than syntactic variation. Another is the development of learning algorithms whose updates are projected onto constraint-preserving subspaces, reducing the burden of post hoc filtering.

A third avenue involves comparative analysis of validation regimes across domains, identifying where consequence-based evaluation suffices and where constraint-indexed approaches are required. Such investigations would transform the present argument from philosophical critique into an empirically testable research program.

\section{Conclusion}

Consequence-Based Utility demonstrates that useful evaluation signals can be extracted even when ground truth is unavailable. Its success highlights the power of consequence as an organizing principle within formal reasoning systems.

Yet consequence alone cannot bear the full weight of validation in open, embodied, and temporally evolving domains. Correctness in such systems emerges not from generativity but from sustained compatibility with constraint.

A comprehensive theory of evaluation must therefore integrate both dimensions. Generativity provides exploration. Constraint provides anchoring. Scientific knowledge advances only where the two remain in tension.

\newpage
\begin{thebibliography}{99}

\bibitem{son2026}
J. Son, A. Radhakrishnan, M. S. Alberti, et al.
\newblock Consequence-Based Utility: Oracle-Free Evaluation of Mathematical Reasoning.
\newblock arXiv preprint arXiv:2602.06291, 2026.

\bibitem{lakatos1976}
I. Lakatos.
\newblock \emph{Proofs and Refutations: The Logic of Mathematical Discovery}.
\newblock Cambridge University Press, 1976.

\bibitem{kuhn1962}
T. S. Kuhn.
\newblock \emph{The Structure of Scientific Revolutions}.
\newblock University of Chicago Press, 1962.

\bibitem{popper1959}
K. Popper.
\newblock \emph{The Logic of Scientific Discovery}.
\newblock Hutchinson, 1959.

\bibitem{polanyi1966}
M. Polanyi.
\newblock \emph{The Tacit Dimension}.
\newblock Doubleday, 1966.

\bibitem{collins1992}
H. M. Collins.
\newblock \emph{Changing Order: Replication and Induction in Scientific Practice}.
\newblock University of Chicago Press, 1992.

\bibitem{shapin1994}
S. Shapin.
\newblock \emph{A Social History of Truth: Civility and Science in Seventeenth-Century England}.
\newblock University of Chicago Press, 1994.

\bibitem{latour1987}
B. Latour.
\newblock \emph{Science in Action: How to Follow Scientists and Engineers Through Society}.
\newblock Harvard University Press, 1987.

\bibitem{boyd2010}
d. boyd.
\newblock Social Network Sites as Networked Publics.
\newblock In Z. Papacharissi (ed.), \emph{Networked Self}. Routledge, 2010.

\bibitem{lessig1999}
L. Lessig.
\newblock \emph{Code and Other Laws of Cyberspace}.
\newblock Basic Books, 1999.

\bibitem{raymond1999}
E. S. Raymond.
\newblock \emph{The Cathedral and the Bazaar}.
\newblock O'Reilly Media, 1999.

\bibitem{merton1973}
R. K. Merton.
\newblock \emph{The Sociology of Science: Theoretical and Empirical Investigations}.
\newblock University of Chicago Press, 1973.

\bibitem{hacking1983}
I. Hacking.
\newblock \emph{Representing and Intervening}.
\newblock Cambridge University Press, 1983.

\bibitem{cartwright1983}
N. Cartwright.
\newblock \emph{How the Laws of Physics Lie}.
\newblock Oxford University Press, 1983.

\bibitem{wilson1998}
E. O. Wilson.
\newblock \emph{Consilience: The Unity of Knowledge}.
\newblock Knopf, 1998.

\bibitem{anderson1972}
P. W. Anderson.
\newblock More Is Different.
\newblock \emph{Science}, 177(4047):393--396, 1972.

\bibitem{smolin2006}
L. Smolin.
\newblock \emph{The Trouble with Physics}.
\newblock Houghton Mifflin, 2006.

\bibitem{taleb2012}
N. N. Taleb.
\newblock \emph{Antifragile: Things That Gain from Disorder}.
\newblock Random House, 2012.

\bibitem{tetlock2005}
P. E. Tetlock.
\newblock \emph{Expert Political Judgment}.
\newblock Princeton University Press, 2005.

\bibitem{crawford2009}
M. B. Crawford.
\newblock \emph{Shop Class as Soulcraft}.
\newblock Penguin, 2009.

\bibitem{sennett2008}
R. Sennett.
\newblock \emph{The Craftsman}.
\newblock Yale University Press, 2008.

\bibitem{scott1998}
J. C. Scott.
\newblock \emph{Seeing Like a State}.
\newblock Yale University Press, 1998.

\bibitem{pariser2011}
E. Pariser.
\newblock \emph{The Filter Bubble}.
\newblock Penguin Press, 2011.

\bibitem{wu2016}
T. Wu.
\newblock \emph{The Attention Merchants}.
\newblock Knopf, 2016.

\bibitem{marres2017}
N. Marres.
\newblock \emph{Digital Sociology}.
\newblock Polity Press, 2017.

\bibitem{goodhart1984}
C. A. E. Goodhart.
\newblock Problems of Monetary Management: The U.K. Experience.
\newblock In \emph{Monetary Theory and Practice}. Macmillan, 1984.

\bibitem{bishop2006}
C. M. Bishop.
\newblock \emph{Pattern Recognition and Machine Learning}.
\newblock Springer, 2006.

\bibitem{amari2016}
S. Amari.
\newblock \emph{Information Geometry and Its Applications}.
\newblock Springer, 2016.

\bibitem{pearl2009}
J. Pearl.
\newblock \emph{Causality: Models, Reasoning, and Inference}.
\newblock Cambridge University Press, 2009.

\bibitem{cover2006}
T. Cover and J. Thomas.
\newblock \emph{Elements of Information Theory}.
\newblock Wiley, 2006.

\end{thebibliography}


\end{document}
